\section{Discussion}
\label{ch:discussion}

This section answers each research question against the evidence of Section~\ref{ch:results}, sets out the mechanism the findings share, relates them to the prior work of Section~\ref{ch:related}, and states the scope and the implications.

\subsection{Answers to the Research Questions}
\label{sec:disc-rqs}

Each research question is answered below against the section of Section~\ref{ch:results} that establishes the answer.

\textbf{RQ1} asked whether the frozen likelihood can be sampled effectively with faithful Langevin dynamics on the discrete manifold, and if not, why. It is answered in the negative across three sections. Section~\ref{sec:results-stepsize} shows that the sampler cannot even be made to move without model-specific calibration, which already denies the strong plug-and-play reading. Section~\ref{sec:results-fallacy} shows, through the direction-versus-magnitude ablation, that in the tested token-substitution settings the input-embedding gradient of frozen autoregressive sequence likelihood provided no reliable proposal advantage over a norm-matched random direction. Section~\ref{sec:results-linradius} gives the mechanism: the surrogate the sampler ranks by is uncorrelated with the true energy change over the whole admissible range of substitution distances, so the gradient is inapplicable rather than merely imprecise, and the input-embedding gradient is in addition structurally blind to the self-term of the energy change. Section~\ref{sec:results-lasttoken} makes the same point exactly at one position: the gradient of the sequence log-likelihood with respect to the final token's input embedding, under the shifted causal indexing in which position $t$ predicts token $t+1$, is exactly zero, while the energy difference between candidate final tokens is exactly the model's conditional log-ratio. Section~\ref{sec:results-diffusion} adds the positive control, in which a score-trained model on the same tokenizer supplies a usable local direction and performs the recovery the gradient-guided sampler did not.

\textbf{RQ2} asked about the role of the Metropolis--Hastings correction in each sampler. It is answered in Sections~\ref{sec:results-quench} and~\ref{sec:results-mh}, and the answer is that the correction plays opposite roles in the two samplers, which is itself informative. For the discrete sampler the correction is beneficial: Section~\ref{sec:results-quench} shows that it removes the quenching artifact and produces smooth early convergence, so here theoretical correctness and empirical performance align. For the continuous sampler the same correction is disabling: Section~\ref{sec:results-mh} shows that it rejects almost every boundary-crossing move, and the decomposition of the acceptance ratio attributes the collapse to the proposal term. The mechanism is that the drift derived from the differentiable pathway of Section~\ref{sec:bg-samplers} is not Lipschitz across a cell boundary; the projected energy of \eqref{eq:bg-cls-energy} is piecewise constant and is never differentiated. The correction does not make the continuous sampler work: it shows that the continuous sampler was functioning as a biased optimizer, so that its apparent successes are successes of optimization under early stopping rather than of sampling.

\textbf{RQ3} asked whether amortization with a GFlowNet escapes the failure. It is answered in Section~\ref{sec:results-gfn}, and the answer is no for the variants tested. Three failure modes appear, capacity collapse, length collapse, and reward hacking, and the unifying experiment shows that Langevin sampling on the amortized energy is indistinguishable from sampling on the untuned base even though the tuning moved that energy by tens of nats. The defensible statement of the finding is that amortization did not escape the observed problems in the tested setup, and that modifying the model through the GFlowNet objective did not produce an energy whose local gradient became more useful. It is not established that no amortized method could work with such a reward, since a GFlowNet's failure can also turn on reward design, training stability, policy parameterization, capacity, exploration, termination handling, or reward scaling.
% SCOPING (Phase 8, evaluation section 3 RQ3: "RQ3 is answered empirically, but the explanation
% should be phrased as: Amortization did not escape the observed problems in the tested setup,
% and modifying the model through the GFlowNet objective did not produce an energy whose local
% gradient became more useful. That is more defensible than implying a general failure of
% amortized inference.") The evaluation's formulation is used here verbatim in substance.

\textbf{RQ4} asked whether an additive constraint can steer generation on this energy landscape. The answer has two statements, which are kept apart because they concern two different landscapes.

First: on the frozen autoregressive energy landscape, plug-and-play constraint steering failed. Section~\ref{sec:results-constrained} reports that on the discrete sampler the trustworthy paired contrast between a real and a randomized constraint direction is essentially zero, and that the one setting in which a constraint-direction signal does appear, the constraint-only arm on a continuous baseline, is entangled with the departure from the fluent manifold that a constraint-only objective causes, so a classifier scoring that text is scoring a distribution it was not trained for.

Second: on a score-trained diffusion landscape, constrained steering became partly possible, subject to fluency and classifier-alignment limits. Section~\ref{sec:results-diffusion-guided} reports that a classifier trained on the noisy states it scores moves an independent judge, in one direction per setting, once a trust region holds fluency fixed, with the reachable direction tracking how far the guiding and judging instruments agree on the text being scored.

These are not one answer. The second uses a different, score-trained landscape and is therefore partly a repair experiment rather than a direct answer about the original autoregressive energy, and it is labelled exploratory where it is reported.
% SPLIT (Phase 8, evaluation section 3 RQ4: "The thesis should separate two conclusions more
% sharply: On the frozen autoregressive energy landscape, plug-and-play constraint steering
% failed. On a score-trained diffusion landscape, constrained steering became partly possible,
% subject to fluency and classifier-alignment limitations. The material is there, but the
% current thesis risks blending those statements.") The two statements are now separate
% paragraphs and are not blended anywhere in the thesis.

\subsection{The Unified Mechanism}
\label{sec:disc-unified}

The failures catalogued above share one property of the object all the methods use. Teacher forcing optimizes each local conditional $p_{\text{LM}}(x_t \mid x_{<t})$ to predict the next token given a correct prefix drawn from the data, and modern models satisfy that objective superbly. What the objective never does is constrain the behaviour of the joint log-likelihood $\log p_{\text{LM}}(\bm{x})$ as a function of $\bm{x}$ when $\bm{x}$ is varied off the data manifold, which is the regime in which every energy-based sampler operates. The joint likelihood is never asked to be smooth, never asked to be a calibrated measure of quality, and never asked to have a gradient that points toward better sequences. The results show that it has none of the three: not smooth, per the projected landscape of Section~\ref{sec:results-mh}; not a good measure of quality, per the likelihood trap of Section~\ref{sec:results-trap}; and not usefully differentiable, per Sections~\ref{sec:results-fallacy} and~\ref{sec:results-linradius}.

This thesis refers to the result of Section~\ref{sec:results-fallacy} by the shorthand \emph{gradient fallacy}, and the shorthand is used only here, with its scope attached: what was measured is that in the tested token-substitution settings, the input-embedding gradient of frozen autoregressive sequence likelihood provided no reliable proposal advantage over a norm-matched random direction. The word does not assert a conceptual error across the field, and it names a bounded empirical finding rather than a general property of autoregressive likelihoods.
% RETAINED ONCE (Phase 8, evaluation section 1: "The term 'gradient fallacy' is particularly
% risky. The experiments support the claim that the tested autoregressive gradient is not useful
% in the investigated sampling settings. Calling it a 'fallacy' suggests a broader conceptual
% error across the field. That may be defensible in the discussion, but it is rhetorically
% stronger than the evidence warrants as a recurring results-section label. I would retain
% perhaps two memorable labels, for example 'linearization failure' and 'likelihood trap', but
% use more neutral section headings elsewhere.") The term is now absent from every heading and
% from every other paragraph in the thesis, and appears only in this one, immediately scoped.
% AUTHOR DECISION: this is the one default the author may wish to overrule. Deleting this
% paragraph removes the term from the thesis entirely, and nothing else depends on it.

There is a mechanical reason the \emph{gradient} in particular is the defective part, and it follows from the decomposition of Section~\ref{sec:results-linradius}. The only term in the likelihood that registers how well a token fits its own left context enters through a discrete index into the output softmax, not through the continuous input embedding the sampler differentiates. The input-embedding gradient can therefore register a substituted token only through that token's effect on the \emph{later} tokens it feeds forward. That blindness is partial in the interior of a sequence and total at the final position, where the gradient of the sequence log-likelihood with respect to the last token's input embedding, under the shifted causal indexing, is exactly zero while the energy difference between candidate final tokens is exactly the model's conditional log-ratio. What fails is not the information in the energy but the attempt to extract a search direction from its input-embedding gradient, and the gradient-free baselines of Section~\ref{sec:results-baselines} confirm the split directly: a single top-$k$ rescoring pass and a Gibbs sampler both recover coherent text on the energy the Langevin samplers cannot navigate.

Two considerations bear on how far this account can be pushed. The first is that it is reached by two routes sharing no machinery. The gradient-guided route differentiates the likelihood and fails, because its surrogate is uncorrelated with the true energy change and the correction, applied honestly, rejects the moves that change a token. The amortized route never touches the gradient, treating the likelihood only as a scalar reward, and it fails too, the learned policy collapsing in whichever direction the reward's implicit incentives push it. The unifying experiment closes the gap between them: running the gradient-guided sampler on the amortized energy reproduces the behaviour of running it on the base. The combined GFlowNet and diffusion results support the interpretation that the training objective, rather than only the sampling algorithm, is a key source of the failure.
% SCOPING (Phase 8, evaluation section 8 formulation 3: "The evidence supports this as a
% plausible mechanism, especially with the diffusion control, but 'because' is stronger than
% 'consistent with'."). Prior wording: "the pathology survives amortization because it is a
% property of the training objective", and "The energy, and not the sampler, and not the search
% procedure, is the problem."

The second is that the same account, read constructively, predicts what would restore the missing direction, and Section~\ref{sec:results-diffusion} bears that prediction out. A model trained by score matching rather than by teacher forcing supplies the direction the autoregressive gradient lacked, the linearization correlation moving from statistically zero to clearly positive; that direction works inside this thesis's own exact-energy Metropolis chain, where swapping only the proposal lifts recovery from zero to thirty-nine percent. The capability the plug-and-play promise assumed to be free is available at the price of training a model to have it.

The likelihood trap qualifies the picture in one further respect. Because low energy is degenerate text on these models, the objective both routes pursue is not the objective one would want. The inability of the gradient-guided samplers to reach low energy is therefore not straightforwardly a loss: the runs that do optimize the energy effectively, beam search on the length-normalized GFlowNet being the clearest, are the runs whose output degenerates into the strings displayed in Section~\ref{sec:results-trap}.
% COMPRESSED (Phase 8, evaluation section 2.E: "Section 6.2 again unifies all failures under
% one defect [...] This is the biggest macro-level repetition."). Five long paragraphs reduced
% to five short ones; the last-token argument, previously stated at full length in both
% Section 5.12 and here, is stated once here in summary and once there in full. Removed
% theatrical phrasings, preserved: "the strange silver lining", "protected from their own
% objective by their inability to optimize it", "This is the deepest sense in which the frozen
% likelihood is not an energy worth sampling", "Two methods with nothing in common beyond the
% energy they use should not fail in ways that trace back to the same property of that energy,
% unless the energy is indeed the common cause", and the closing "Making this last-position
% argument exact [...] is the natural culmination of the diagnosis and the immediate next
% analysis the work points to" (which was written before the analysis was run).

\subsection{Connections to the Related Work}
\label{sec:disc-related}

The likelihood trap measured in Section~\ref{sec:results-trap} is a direct, quantitative reproduction on the study's own models of the neural-text-degeneration phenomenon that \citet{holtzman2020curious} identified: that the most probable sequences under a neural language model are degenerate, and that maximization-based decoding produces worse text than sampling. Section~\ref{ch:related} presented that as a known hazard of decoding. The contribution here is to re-derive it as an obstacle for energy-based sampling specifically, since in that framing the degeneracy of high-probability text says that the \emph{target} of the sampling procedure is undesirable rather than that a decoding heuristic should be avoided.

The anisotropy results of Section~\ref{sec:results-aniso} confirm, on GPT-2 Large and Llama-3, the representation-degeneration and embedding-anisotropy phenomena documented by \citet{gao2019representation} and \citet{ethayarajh2019contextual}, and add a consequence those works did not draw: because the anisotropy scale differs by roughly a factor of three between the two models, a single step-size schedule cannot transfer between them, which makes the samplers model-specific in a way that qualifies their claim to be plug-and-play.

The Metropolis--Hastings breakdown of Section~\ref{sec:results-mh} is the empirical counterpart to a gap noted in Section~\ref{sec:rel-energy}: that the energy-based decoding literature frequently omits the correction, and that where it is included, the difficulty of computing the reverse proposal on a projected landscape is seldom examined. This thesis implements the correction faithfully, following the convergence theory of \citet{roberts1996exponential}, and measures what the common omission conceals. On that reading, methods such as COLD \citep{qin2022cold} and MuCoLa \citep{kumar2022gradient}, which the continuous sampler of this thesis follows, are better understood as biased optimizers than as samplers from the stated energy, and their reported successes as successes of optimization under early stopping. This is a clarification of what the methods do rather than a criticism of those works, which are careful about what they claim. The same reading applies to the annealed stochastic-gradient Langevin dynamics of \citet{welling2011bayesian}: without the correction, the annealed schedule its guarantee relies on becomes the quenching effect of Section~\ref{sec:results-quench}, freezing the chain into whichever mode it last occupied.

The GFlowNet results of Section~\ref{sec:results-gfn} qualify, without contradicting, the picture presented by \citet{hu2024amortizing}, whose formulation and codebase this thesis builds on. Their work showed that amortized sampling from a frozen-model posterior can yield diverse, transferable samples where the reward is well-behaved, and nothing here disputes that. The narrower finding is that when the reward is the raw frozen likelihood used as an energy over free-form sequences, the learned policy collapses in the ways Section~\ref{sec:results-gfn} documents. The axis that reconciles the two is whether the reward defines a landscape with a coherent-text optimum.

The gradient-free samplers isolated in Section~\ref{sec:rel-energy}, Mix-and-Match \citep{mireshghallah2022mixmatch} and the twisted sequential Monte Carlo of \citet{lew2023sequential} and \citet{zhao2024probabilistic}, confirm rather than contradict the central result. Each samples from a frozen-model energy without differentiating it, and each works. The Gibbs baseline of Section~\ref{sec:results-baselines} makes the point in-platform, on the identical energy the Langevin methods use. The frozen likelihood is therefore a serviceable object to \emph{evaluate}, and to search without gradients, and an unusable one to \emph{differentiate} in the settings tested here, which is why the thesis narrows its claim to the gradient rather than the energy.

\subsection{Scope and Limitations}
\label{sec:disc-scope}

The thesis shows that in the settings tested, the frozen likelihood of a standard autoregressive language model is a poor energy for gradient-guided and amortized sampling on discrete text, and it explains why in terms of the training objective and the geometry of the token embedding space. It does not show that energy-based controllable generation is impossible in general.

One formulation deserves particular care, because the thesis set out to test it. The experiments refute the premise that the frozen autoregressive sequence-likelihood gradient can supply the required local score for this class of Langevin-based methods without additional training. They do not refute the training-free premise as such. Training-free methods that do not differentiate the raw joint likelihood remained viable in this study's own results: the gradient-free Gibbs sampler of Section~\ref{sec:results-baselines} matches the best gradient-guided Langevin configuration on the identical energy, and a single top-$k$ rescoring pass beats every configuration in the grid, both without any training and without any gradient. Methods that use lookahead, conditional resampling, or another representation may likewise remain viable, and Section~\ref{sec:rel-energy} reviews published members of that family.
% SCOPING (Phase 8, evaluation section 8 formulation 4: "'The training-free premise is refuted'
% [...] It is refuted for the tested framework, not necessarily for every possible training-free
% control method built on autoregressive models. Methods that do not differentiate raw joint
% likelihood, use lookahead, use conditional resampling, or operate in another representation
% may remain viable."). The evaluation's formulation is used verbatim, and the note that
% gradient-free training-free methods remained viable is stated explicitly, on this thesis's own
% results. Prior wording, in the abstract and the conclusion: "the training-free premise on which
% the whole framework rests is tested and refuted".

Several further caveats bound the empirical claims. The KL-divergence metric depends on the position of the corrupted token within the sequence, which is why, as Section~\ref{sec:meth-metrics} explained, the central results are argued from the absence of a gap between conditions rather than from small measured differences; the gradient result should be read as ``no reliable difference'' rather than as an exact numerical equality. The Llama-3 results support only qualitative cross-architecture claims, because the tokenizer and embedding scale differ from the GPT-2 family. The linearization result is that the surrogate is uniformly uninformative across the admissible range of distances rather than informative up to a sharp radius and useless beyond it, so the phrase ``linearization radius'' summarizes a decay rather than naming a threshold. The GFlowNet result covers the variants trained here and not amortized inference in general, as Section~\ref{sec:disc-rqs} states. And the constrained-generation extension establishes the absence of reliable steering on this energy, not the impossibility of steering on a better one.

Finally, the relationship between the title of this thesis and its content. The original objective was controllable generation via faithful Langevin dynamics, planned in two stages: establish that energy-guided sampling on a frozen model is viable, then add the constraint term and perform sentiment and toxicity control. The first stage did not hold, and the work turned to the question of why rather than building control on a foundation it had shown to be unsound. What it delivers is a diagnosis of the failure of a prerequisite, with the mechanisms identified, measured, and cross-checked, and with the constrained-generation extension included as evidence of the downstream consequence, rather than a demonstration of successful control.

\subsection{Implications}
\label{sec:disc-future}

Locating the defect in the training objective rather than in the sampler makes a testable prediction: a model whose objective does shape a score function over sequences should not fail in the same way. Diffusion language models are such models \citep{li2022diffusion, lou2024discrete}, because denoising is equivalent to learning the score \citep{vincent2011connection, song2019generative, ho2020denoising}. Section~\ref{sec:results-diffusion} runs that control on the score-entropy discrete diffusion model of \citet{lou2024discrete} at two scales, and every arm of the prediction is borne out: the linearization correlation is clearly positive where the autoregressive gradient was statistically zero, the model natively recovers the flipped token, and replacing only the proposal direction inside this thesis's exact-energy Metropolis chain lifts exact recovery from zero to thirty-nine percent. The correlations are modest and the diffusion models are smaller and differently trained than the frozen backbone, so these are pilots and are read as such.
% COMPRESSED (Phase 8, evaluation section 2.E: "Section 6.5 again states implications and the
% decisive positive control", and section 5 item 10). Section 6.5 previously re-narrated the
% design and outcome of the diffusion control at length, in a passage largely written before the
% control was run and then amended in place. The outcome is reported once, in Section 5.13, and
% summarized in one paragraph here. The Phase 3 "designed but not yet run" wording and its
% Phase 4 replacement are preserved below for provenance.
% PRIOR WORDING (Phase 3, before the control was run):
% "The concrete instrument for the control is the score-entropy discrete diffusion model of
% \citet{lou2024discrete} ... It is honest to report the status of this experiment precisely,
% because it is designed but not yet run. ... running it is the most valuable single step this
% thesis leaves open."
% PRIOR WORDING (Phase 4, replaced in Phase 8): the section was titled "Implications and the
% Falsifying Experiment" and restated the full linearization, recovery, hybrid and guidance
% outcomes, plus the reframing of the promise, in two long paragraphs.

Two implications follow for practice. The first concerns what inference-time control on a frozen autoregressive model can be built from. The energy is usable and its input-embedding gradient, in these settings, is not, so a method that evaluates or reranks under the frozen likelihood has a foundation and one that differentiates it does not, at a cost that Appendix~\ref{app:cost} shows also favours the former. The second concerns what would have to change to recover the gradient-guided route. If the problem is that the likelihood is not a score function, one option is to make it one, by fine-tuning the model so that its sequence log-density is aligned with a target energy under a score-matching objective. That is a lighter intervention than reinforcement-learning-style training of a GFlowNet and is directed at the identified deficiency rather than at a downstream symptom. Only once an energy landscape is known to be navigable does the additive constraint term of the plug-and-play framework become meaningful, since only then is there a fluency energy a sampler can follow for the constraint to be added to.
